\chapter{A Weekly Literature Review}
\label{ch:weeklyliterature}

The previous chapter ran one agent session and checked its work. This chapter
builds something that runs without you, every Monday, and whose output you can
trust for the same reason: every claim in it can be checked. The task is the
one \Cref{ch:keepingup} calls impossible to do completely, keeping up with the
literature. A weekly review does not make it possible. It makes the part that
machines do well, fetching, merging and verifying, routine and recorded, and
leaves the reading to you.

The design has two layers that are never mixed (\Cref{fig:twolayers}). The
first is a set of R scripts with fixed queries, written once and run by a
clock. They fetch the week's new items from five sources, merge duplicates,
resolve every identifier, and write a Markdown digest, a BibTeX file and a log
of every query. No model is involved. The second layer is an assistant, a
Claude Code skill built on the rules KYBELE uses (\Cref{ch:localfirst}). It
reads the week's files, searches Consensus, and writes notes in a separate
file that says a model wrote it. It can add a reference to a bibliography only
through a script that resolved the reference first. The reason for the split
is the argument of \Cref{ch:reproducibility}: code that a model wrote can be
kept and rerun, but an answer a model gave cannot.

\begin{figure}[htbp]
\centering
\fig{ch11-two-layers}
\caption{The two layers. Scripts in R run on a clock and produce the record;
the assistant reads that record and may write only its own notes. The one
arrow from the assistant into the bibliography passes through a script that
resolves the reference.}
\label{fig:twolayers}
\end{figure}

By the end you should have a repository on GitHub that fills itself each week
and that anyone can read, and a file that states your interests as queries.
You should also have a way of letting a model summarise what came in without
letting it invent anything. Allow about three hours to set it up and half an hour each Monday to
read the result. You need R, Git and a GitHub account; the assistant needs
Claude Code and a Consensus account. The \code{workshop/literature} folder
holds a frozen week, so that the exercise runs offline and its counts can be
compared with the ones printed here. The repository this chapter describes is
listed in \Cref{app:tutorials}, and the versions used are in
\Cref{app:stateofplay}.

\section{The sources}

Five services offer a search interface that a script can query by date, and
each covers a different part of the literature. Europe PMC indexes the life
sciences with abstracts and includes preprints; a query can ask for items
first published between two dates. OpenAlex covers every field and most
publishers. It returns abstracts as an inverted index that the script
rebuilds, and it stems words. The acronym MIxS therefore matches every paper
that mentions a mix. Computer science appears first on arXiv, including most
of the work on language models that this book cites; its API takes a date
range and a category. The bioRxiv API returns everything in a category for a
date range, thirty items at a time, and the script keeps what matches a topic's
terms. Biodiversity PMC \parencite{pasche2026biodivpmc} is the
biodiversity-specific index built on the SIB Literature Services, which KYBELE
also uses. Its \code{medline} collection is current. Its full-text \code{pmc}
collection ran about two months behind for the author's queries when this was
tested. The Plazi treatments it holds \parencite{agosti2022treatments} carry
no publication date in the index, so they cannot be searched by week.

Consensus is the sixth source and the odd one. It searches a large index of
papers with a ranked, semantic query, which suits questions rather than terms.
It has no public API, only an MCP server (\Cref{ch:aireadydata}) that a
harness can call. It therefore belongs to the assistant layer, and the scripts
never depend on it.

Two properties of these sources shape the design. They index with a delay,
and not the same delay, so an item can appear in one source this week and in
another next week. And a preprint and its journal version arrive weeks apart
under different identifiers. The scripts therefore fetch some days before the
week as well, remember everything they have ever listed, and merge by
identifier before they merge by title.

\section{Step 1: your interests, as queries}

KYBELE's first rule is that the vocabulary lives in data, looked up, never
generated by the model. The weekly review keeps the same rule. A file,
\code{config/interests.yml}, holds the topics, and for each topic the terms
that mark it, one query per source, and the questions the assistant may put
to Consensus. You write it by hand. It is the one part of the system that no
prompt in this chapter asks an agent for. One topic looks like this:

\begin{lstlisting}[caption={One topic in config/interests.yml}]
topics:
  - id: collembola
    label: Collembola, soil fauna and KYBELE
    terms: [Collembola, collembolan, collembolans, springtail, springtails,
            "Fauna Europaea", "taxonomic treatment"]
    europepmc: 'TITLE_ABS:(Collembola OR springtail*)'
    openalex: 'Collembola OR springtail OR springtails'
    biodivpmc: 'Collembola OR springtail OR springtails'
    biodivpmc_cols: [medline, pmc]
    biorxiv_categories: [ecology, evolutionary_biology, zoology]
    arxiv: null
    consensus:
      - "Collembola responses to soil disturbance"
      - "springtail new species description"
\end{lstlisting}

Fourteen topics took an afternoon to write, from this book, the author's notes
and the projects of \Cref{ch:localfirst}. The counts of the first weeks then
show which queries are too broad and which find nothing, and the file is tuned
from them. It is under version control, so the log of every week records
which version of the vocabulary produced it.

\section{Step 2: fetch, and keep what came back}

Set up the repository as in \Cref{ch:firstsession}: a folder, \code{git init},
a \code{README.md} that says what the system is for, and an \code{AGENTS.md}
with the rules. Then ask for the first script:

\begin{lstlisting}[caption={The first prompt}]
Read README.md and AGENTS.md. Write R/00_common.R and R/01_fetch.R in base R
plus httr2, jsonlite, yaml and xml2; no other packages. 01_fetch.R takes an
ISO week label such as 2026-W42, or defaults to the week that ended
yesterday, reads config/interests.yml, and for every topic queries Europe
PMC, OpenAlex, arXiv, bioRxiv and Biodiversity PMC for items first published
in that week plus the overlap days in the config. Save every response
unchanged under data/raw/<week>/<source>__<topic>__<page>.json, or .xml for
arXiv. Send a User-Agent with the contact address from the config. Wait
three seconds between arXiv calls and one second otherwise. Retry a failed
request four times with growing delays. Never stop the run because one
source failed: record the failure and continue. Print one line per source
and topic with query, pages, items and seconds. Done when
Rscript R/01_fetch.R 2026-W42 writes the files and the number of items
counted from the files with jsonlite equals the printed totals.
\end{lstlisting}

Three things in this prompt carry the method. Every response is saved
unchanged, before anything is parsed, so that a later question about what a
source returned has an answer. The request identifies you, which the
services ask for. And the condition at the end is one that can fail: the
counts the script prints must equal the counts in the files it wrote.

This was run for this book on 14 October 2026, for the week of 5 to 11
October. The fourteen topics and the people queries made 51 requests to the
five sources, with their pages, in about three and a half minutes. They
reported 2,482 items: 1,331 from bioRxiv's six categories, 538 from arXiv,
492 from OpenAlex, 92 from Europe PMC and 29 from Biodiversity PMC.
Seventeen queries returned nothing, and the review says which. The number is
large because the first run of a new week fetches the ten overlap days too.
It is also large because the terms, not the fetch, decide what bioRxiv
contributes: 1,242 of its records matched no topic and were set aside. An
earlier attempt had
been smaller and wrong in ways that only the counts revealed. Of the 156
OpenAlex hits for Collembola, 149 were Catalogue of Life records, until the
query asked for articles, preprints and reviews only. The acronym MIxS
brought 1,599 items about mixing. The arXiv query for agents returned 424 in
one week, because the word harness, as a verb, matched 128 of them. Each fix
is a line in the configuration, with the date and the reason.

\section{Step 3: one table, no duplicates}

\begin{lstlisting}[caption={The second prompt}]
Write R/02_normalise.R. It reads every file under data/raw/<week>/ and
writes data/normalised/<week>.csv with one row per distinct item and the
columns listed in README.md. Merge duplicates first by DOI, lower case and
without the https://doi.org/ prefix, then by arXiv identifier, treating an
OpenAlex DOI of the form 10.48550/arXiv.ID as that identifier, then by
normalised title: lower case, ASCII, letters and digits only. When merging,
keep the union of topics and sources, the journal venue over the preprint
venue, and the longest abstract, cut to sixty words. Print rows read and
rows after each of the three merges. Done when
Rscript R/02_normalise.R 2026-W00 prints 6 read, 4 after DOI, 4 after
identifier, 3 after title, and the row for DOI 10.1234/fixture-1 lists two
sources.
\end{lstlisting}

\begin{figure}[htbp]
\centering
\fig{ch11-week-pipeline}
\caption{One week through the scripts. Five sources are saved as received,
merged into one table, and checked; an item whose identifier does not resolve
is listed as not verified and gets no bibliography entry.}
\label{fig:weekpipeline}
\end{figure}

The merge order matters (\Cref{fig:weekpipeline}). A DOI is an identity, an
arXiv identifier is one too, and OpenAlex records arXiv preprints under a DOI
that contains the identifier. A title is not an identity. Two records with the
same normalised title are probably the same work, but the fixture week in the
repository contains a pair that is not, and the check in the next step catches
it. The fixture is six invented records across four sources, and the prompt's
last sentence states what the script must print for them. Week \code{2026-W00}
does not exist in any calendar, which is the point: it can never be confused
with a real week.

\section{Step 4: resolve every identifier before it enters the bibliography}

\Cref{ch:scientificworkflow} resolves every reference before it enters a
manuscript, because about 147,000 invented citations appeared in 2025 alone
\parencite{zhao2026citations}. The weekly review applies the same rule one
step earlier: nothing enters a \code{.bib} file that a registrar has not
confirmed. The prompt points the agent at the book's own script and asks for
the same check in R:

\begin{lstlisting}[caption={The third prompt}]
Read scripts/verify_bib.py from the book repository, then write
R/03_check.R to do the same for data/normalised/<week>.csv. For each DOI
fetch https://api.crossref.org/works/<doi>, and if Crossref answers 404 fetch
https://api.datacite.org/dois/<doi>, where Zenodo and Figshare register.
Compare the returned title with ours after normalising, using 1 minus the
edit distance over the longer length, threshold 0.85; also accept the pair
when one title is the other without its subtitle. Set verified TRUE or
FALSE with a reason. Check arXiv identifiers through the arXiv API id_list,
at most 100 per request, three seconds apart. Save the responses under
data/raw/<week>/crossref.json and arxiv-ids.xml (one file per batch) and
read them instead of the network when --offline is given. Skip identifiers
already in data/seen.csv: they were verified when first seen. Write
data/normalised/<week>-checked.csv. Add --self-test with a good DOI, a DOI
that does not exist, a real DOI with a deliberately wrong title, a good
arXiv identifier and one that does not exist; it passes only if the three
bad cases fail and the two good ones pass. Done when
Rscript R/03_check.R --self-test prints SELF-TEST PASSED and names the three
failures.
\end{lstlisting}

The self-test is the part to insist on. A verifier that cannot fail is worse
than none, because it reports success for everything. Planting a DOI that does
not exist, a real DOI with the wrong title and an arXiv identifier that does
not exist, and requiring all three to fail, is a check that the check works.

The first live week taught the verifier three things the fixture had not.
The arXiv API refused 480 identifiers in one request, so the check now sends
a hundred at a time. Of the 53 DOIs Crossref did not know, 41 were Zenodo
records, which DataCite registers, so the check now asks DataCite when
Crossref has never heard of a DOI. In the final run DataCite confirmed 105
items. And two of the four title mismatches were a subtitle that one source
had kept and another dropped, so a shorter title that begins the longer one
now passes, with the reason recorded. In the final run 540 DOIs and 553
arXiv identifiers were checked in a little over five minutes: 1,071 items
passed and 35 did not. Fifteen DOIs resolved at neither registrar, thirteen
items had no identifier at all, and seven had a title that differed from the
registrar's. A DOI can take days to register, which is why the scripts fetch
ten days of overlap and try an unverified item again while it stays in the
window. All 35 are listed in the review as not verified, and none is in the
bibliography.

\section{Step 5: the week's files and the log}

\begin{lstlisting}[caption={The fourth prompt}]
Write R/04_write.R. From data/normalised/<week>-checked.csv and data/seen.csv
it writes reviews/<week>.md and bib/<week>.bib, appends new entries to
bib/all.bib and new rows to data/seen.csv, and regenerates index.md. Only
verified items get a bib entry; the others are listed at the end of the
Markdown under "Not verified" with the reason and never enter a .bib file.
Keys are authorYEARkeyword: first author's surname in ASCII, year, first
title word longer than three letters that is not a stop word, made unique
against all.bib with a letter suffix. Preprints are @online entries with
eprint and eprinttype for arXiv or the bioRxiv DOI, and the Markdown marks
them preprint. Items already in seen.csv are not listed again. A preprint
with a higher version is listed once under "Updated preprints". An article
whose DOI or normalised title matches a seen preprint is listed once under
"Now published" and the seen row gets published_doi. Record every source
and topic that returned nothing under "Searches that returned nothing".
Escape ampersand, percent, underscore, hash and braces. Print items in,
new, seen before and bib entries written. Done when running the script
twice for 2026-W00 leaves all.bib and seen.csv unchanged the second time,
shown by git diff --stat, and biber --tool bib/2026-W00.bib reports no
errors.
\end{lstlisting}

The last condition, that a second run changes nothing, is the test of a script
that writes to files it also reads. A review that listed the same paper twice
because the script ran twice would soon be ignored. The sections that record
absence, ``Not verified'' and ``Searches that returned nothing'', are there
for the same reason as the gold set in \Cref{ch:testing}: an empty result is a
result, and it must be visible. The fifth prompt asks for the script that runs
the other four and writes the log.

\begin{lstlisting}[caption={The fifth prompt}]
Write R/run_week.R. It runs R/01_fetch.R to R/04_write.R in order for one
week, given as an argument or defaulting to the week that ended yesterday,
passes --offline through, and writes provenance/<week>.log.md with: run
time in UTC, the git commit of the config, R and package versions, the md5
sum of every raw file, every query with its URL without any key, pages,
items and seconds, the normalise counts, DOIs checked, passed and failed by
key, files written, every failure and retry, and the total time. Set up
renv with httr2, jsonlite, yaml and xml2 only, and snapshot. Done when
Rscript R/run_week.R 2026-W00 --offline finishes without error, the log
contains every section above, and renv::status() reports the project is
synchronised.
\end{lstlisting}

The log is the instrument reading of the week: which queries ran, what each
returned, which identifiers resolved and which did not, and how long it took.
With the raw responses kept beside it, a week can be rerun and the two runs
compared. That is the standard of \Cref{ch:reproducibility}, applied to a
process that reads from services that change. The same queries run a month
later will not return the same items, and the log is what makes the
difference visible.

Everything the review produces is a plain file: Markdown, BibTeX, CSV, a log.
That is deliberate. Files are what a person, a script and an agent can all
read without an intermediary (\Cref{ch:agentic};
\textcite{piskala2026files}), and GitHub renders the Markdown, so the review
needs no website. The raw responses are the one exception. They hold the full
abstracts, which are the publishers' text, so they stay out of the public
repository and are kept as an artefact of the run, with their checksums in
the log.

\section{Step 6: let a machine run it on Monday}

\begin{lstlisting}[caption={The sixth prompt}]
Write .github/workflows/weekly.yml: on a schedule every Monday at 06:00 UTC
and on workflow_dispatch with an optional week input. On ubuntu-latest:
check out, install R 4.5.1 with r-lib/actions/setup-r using public RSPM
binaries, install libxml2-dev and libcurl4-openssl-dev, restore the renv
library, run Rscript R/03_check.R --self-test, run Rscript R/run_week.R
with the week input if given, upload data/raw/ as an artifact, and commit
reviews/, bib/, data/normalised/, data/seen.csv, provenance/ and index.md
with the message "Week <week>: <n> new items" when there is anything to
commit. OPENALEX_API_KEY comes from a repository secret; the scripts read
it from the environment and never print it. Store no other secret. Done
when a workflow_dispatch run for 2026-W41 is
green, its commit is on main, and the job log shows the self-test passed
before the fetch.
\end{lstlisting}

GitHub Actions is the clock. It runs on GitHub's machines, so your laptop
need not be awake. Its run logs are public, so the record of each week lies
not only in the repository but in the run that produced it. The workflow
stores one secret, an OpenAlex key, which the scripts read from the
environment and never print. The key is not optional there. OpenAlex meters
requests without a key against a free daily budget shared by every machine
behind one address, and a GitHub runner shares its address with many. On the
author's own machine, five full fetches in one afternoon spent the budget.
The service then asked for a retry ten hours later. The self-test runs before
the fetch on every Monday, so a verifier that stopped working would stop the
run rather than pass everything.
GitHub's documentation warns that a public repository's schedule is disabled
after sixty days without activity. A scheduled run can also be delayed when
the service is busy, and schedules run only on the default branch. The weekly
commit is itself activity, and a missed week can be run by hand with the week
as input.

The same scripts run on a Mac under \code{launchd}, the system's scheduler,
and the repository holds the file that does it. That is the route for the
assistant, which needs a Consensus connection that only an interactive login
can create. A third route, a scheduled task in the harness's own cloud, had
not been confirmed to offer R or Consensus when this was written, and the
repository says so.

\section{Step 7: an assistant in KYBELE's shape}

The scripts produce a digest of abstracts. Reading it is your job, but a
model can help: group what came in, draw the connections to your projects,
and put to Consensus the questions that term searches cannot. The danger is
the one \Cref{ch:powers} describes, a fluent summary of papers the model did
not read, with references it did not retrieve. KYBELE's rules are the answer.
The assistant is written as a \term{skill}, a folder of instructions and
scripts that the harness loads for one recurring task
\parencite{anthropic2026skills}. The seventh prompt asks for it.

\begin{lstlisting}[caption={The seventh prompt}]
Read docs/assistant-rules.md; its rules are not to be weakened. Write
.claude/skills/weekly-literature/SKILL.md in the Claude Code skill format
with name, description, argument-hint, allowed-tools and
disable-model-invocation. The skill takes a week label, reads
config/interests.yml, reviews/<week>.md and bib/<week>.bib, runs the
Consensus queries listed in the config for each topic through the Consensus
MCP tool, three at a time, and appends every query with UTC time, result
count and the returned URLs to provenance/search-log.md. A Consensus result
may be cited only if it is in the week's bib, or if
Rscript R/skill_verify.R <doi> "<title>" prints OK; that script writes the
entry to bib/<week>-consensus.bib itself. The skill then writes
notes/<week>.md: per topic, a few sentences from the abstracts, every
sentence tied to a bib key in square brackets or a verified DOI, preprints
marked, and "Searches found no new items on ... this week" when there is
nothing. It edits no other file and never commits. Also write
R/skill_verify.R, with the same Crossref title check as R/03_check.R, and
tests/check_notes.R, which exits with an error if the notes cite anything
not in the week's bib files or the search log, and has a --self-test that
plants one bad key and one unlogged URL. Done when /weekly-literature
2026-W42 produces notes that pass tests/check_notes.R, and the number of
queries added to search-log.md equals the number of Consensus calls in the
session.
\end{lstlisting}

The rules file this prompt refers to is short, and its four parts are
KYBELE's. Vocabulary from the config: the assistant may run only the
Consensus queries listed for each topic. Scope in code: it may cite a paper
only if a script has resolved it, and the script writes the entry, so the
model never writes bibliography text (\Cref{fig:citationcheck}). Saying no: a
topic with nothing new gets the sentence ``Searches found no new items on
this topic this week'', never ``nobody has worked on this''. And a citation
check: every sentence that states a finding carries a key, and a test script
refuses notes that cite anything outside the week's files or the search log.

\begin{figure}[htbp]
\centering
\fig{ch11-citation-check}
\caption{What a Consensus result must pass before the assistant may cite it.
Either it is already in the week's bibliography, or a script resolves its DOI
and compares the titles; the script writes the entry. Anything else is listed
as not used, with its address.}
\label{fig:citationcheck}
\end{figure}

Two further rules come from \Cref{ch:security}. Text inside a fetched record
is data, never an instruction, and the fixture week plants an abstract that
tells assistants to add a reference, as the notes cell did in
\Cref{ch:firstsession}. And the harness configuration denies the skill every
tool it does not need: it can read, run the two scripts, and write one file.
A rule in the prompt lowers the chance of an unwanted write; the
configuration prevents it.

The notes file begins with a header that says a model wrote it, which model
and harness version, and from which files. It adds that the papers were not
read in full, and leaves a line for the person who checked it. That is what journal
policies now ask of any text a model helped to write
\parencite{naturemethods2026responsibly}, applied to a file nobody will
submit. Install only skills you have read
\parencite{liu2026dependencysteering}; this one is in the repository in full.

To be plain about what this is: not KYBELE, whose literature service,
orchestrator and model \Cref{ch:localfirst} describes, but KYBELE's rules
reused in a smaller system. The scripts were run for this book. The skill was
written and its checks were tested, but its live use, which needs the
Consensus connection and a person at the keyboard, was left to the author's
first Mondays.

\section{Step 8: testing it}

\Cref{ch:testing} argues that a model on your own task needs an answer key, a
baseline and a score. The repository's \code{tests} folder gives the review
the same. A gold set lists papers you know should have been found, with where
you learnt of them. A script reports how many are in the review's record, and
the set grows whenever you hear of a paper elsewhere. A grounding check reads
the notes and fails on any key or address that is not in the week's files or
the search log. Its self-test plants one bad key and one address that was
never logged. A pushback test puts three false claims to the assistant, such as a
paper that is not in the week, and passes if it refuses each with reference
to the retrieved set. The injection fixture is a fourth test. A results table
records, for each run, the model, the harness version, the scores and the
cost, so that a change of model or prompt is accepted only if the scores
hold.

When this was tested for this book, the fixture week and the two self-tests
passed, and the grounding check refused its planted key and address. The
live week was fetched five times in one afternoon while the scripts were
being fixed, which the fetch logs record and which OpenAlex eventually
refused. The final files therefore came from the complete fetch of the early
afternoon, kept as received, and from a Biodiversity PMC fetch on its own
after its page-size error was found. Every identifier was then verified
live, and the log says so in its first lines. Biodiversity PMC had failed in
every earlier
run because the request carried its page size twice. Its full-text
collection then returned nothing for the window, since on that day its newest
Collembola record was two months old, while its MEDLINE collection was
current and returned 29 items. The week's review lists 1,071 verified items
under thirteen topics, about two thirds of them preprints, which is far more
than anyone will read. The two largest topics, agents and language models in
biodiversity, hold two thirds of them, and the review caps each topic at 25
listed items and counts the rest. That is the state of a first week: the
vocabulary is too wide, and the counts say where. The assistant's notes and
the pushback test were not run, because they need the Consensus login and
the author at the keyboard. The gold set holds twelve papers from this book's
bibliography, all older than the week, so the recall test has nothing to
score yet.

\section{What to take away}

\begin{itemize}
  \item Two layers, never mixed. Scripts with fixed queries produce the
    record; a model may comment on it, in a file that says so.
  \item A reference enters a bibliography through a script that resolved it,
    never through a model.
  \item Keep what the sources returned and log what was asked. The same query
    will not return the same answer next month.
  \item The vocabulary is yours. Write it by hand, version it, and tune it
    from the counts.
\end{itemize}

\section{Where to go from here}

Add sources: the Biodiversity Heritage Library for historical literature,
Zenodo's BioSysLit community for taxonomic papers, and your institution's
repository. Let a local model write the notes through \code{ellmer}
(\Cref{ch:testing}), under the same rules and the same checks, and compare its
grounding score with the hosted model's. Then build the small project of
\Cref{ch:keepingup} on the weekly record: retrieval over the abstracts you
have accumulated, with citations. \textcite{zhang2025deepresearch} describe
systems that decide for themselves what to search next. This review is the
opposite design, fixed queries with a logged result, which is what makes it
reproducible. The two can meet: a question answered over a record whose every
item was verified before it was stored. Without retrieval, GPT-4o
fabricated citations to recent literature in 78 to 90 per cent of the
OpenScholar tests \parencite{asai2026openscholar}, which is why the notes
are confined to what the scripts fetched. A shared knowledge layer that agents
query, as \textcite{sigillo2026embllibrarian} propose, is where records like
this one could eventually pool.
